\section*{الفصل \ref{ch:b2:muscle-movement} --- التقلص العضلي والوظيفة الحركية}

\begin{solution}{exo:b2:muscle-movement:1}
تنزلق الخيوط الرفيعة على امتداد الغليظة، تجرها رؤوس الميوزين؛
ولا يقصر أي من الخيوط. ويضيق النطاق I والمنطقة H ويتقارب
قرصا Z؛ ويحفظ النطاق A (أي الخيوط الغليظة) طوله.
\end{solution}

\begin{solution}{exo:b2:muscle-movement:2}
(1) ارتباط ATP وانفصال الرأس؛ وحلمهة ATP وزنبرة الرأس. (2) ثم يرتبط الرأس
بالأكتين. (3) ثم يُطلق الفوسفات، فضربة القدرة، ثم يُطلق ADP.
(4) ثم يرتبط ATP جديد فينفصل الرأس. وATP لازم
\emph{لإفلات} الرأس: فدونه يبقى كل رأس مرتبطًا وتنقفل
العضلة --- وهو التصلب.
\end{solution}

\begin{solution}{exo:b2:muscle-movement:3}
دفعة عصبية؛ ثم إطلاق الأستيل كولين وكمون الصفيحة الانتهائية
(\qty{0.6}{ms})؛ ثم كمون فعل عضلي على الليفة ونازلًا في
نبيبات T (\qty{1}{ms})؛ ثم إطلاق الكالسيوم من \omterm{def:b2:cell-differentiation:fibre}{الشبكة الهيولية
العضلية} (\qty{1}{ms})؛ ثم ارتباط الكالسيوم بالتروبونين، وتزحزح التروبوميوزين،
وارتباط الرؤوس (بضع مليثوان)؛ ثم ارتفاع القوة على مدى عشرات
المليثوان.
\end{solution}

\begin{solution}{exo:b2:muscle-movement:4}
\omterm{def:b2:muscle-movement:motorunit}{الوحدة الحركية}: \omterm{def:b2:neurons-synapses:neuron}{عصبون} حركي واحد وكل الليفات التي يعصّبها.
\omterm{def:b2:muscle-movement:motorunit}{والرعشة}: التقلص الوجيز من دفعة واحدة. \omterm{def:b2:muscle-movement:motorunit}{والكزاز}: التقلص
المندمج المستدام من دفعات تصل أسرع مما تخبو
\omterm{def:b2:muscle-movement:motorunit}{الرعشة}. وتُدرَّج القوة بمعدل إطلاق كل وحدة
وبعدد الوحدات المجنَّدة.
\end{solution}

\begin{solution}{exo:b2:muscle-movement:5}
ينزلق كل نصف قسيم \qty{0.2}{\micro m} في \qty{50}{ms}:
أي \qty{4}{\micro m/s}؛ وذلك 20 ضربة قدرها \qty{10}{nm} في
\qty{50}{ms}، أي 400 في الثانية لو كان رأس واحد مرتبطًا طوال الوقت.
\end{solution}

\begin{solution}{exo:b2:muscle-movement:6}
\qty{2}{cm} في \qty{0.1}{s}: أي \qty{0.2}{m/s}، وطولان في الثانية؛ وكل
\omterm{def:b2:cell-differentiation:fibre}{قسيم عضلي} يقصر $2/\num{40000}$ cm $= \qty{0.5}{\micro m}$ في
\qty{0.1}{s}، أي \qty{5}{\micro m/s}.
\end{solution}

\begin{solution}{exo:b2:muscle-movement:7}
$F/F_0 = 0.3125/(v/v_{\max} + 0.25) - 0.25$: فعند 1 و3 و6 أطوال في
الثانية (أي $0.1$ و$0.3$ و$0.6$ من $v_{\max}$): $0.64$ و$0.32$ و$0.12$.
والاستطاعة $Fv$ (بوحدة $F_0\times$ طول في الثانية): $0.64$ و$0.95$ و$0.71$
--- وهي أعظم قرب 3 أطوال في الثانية، أي نحو ثلث $v_{\max}$.
\end{solution}

\begin{solution}{exo:b2:muscle-movement:8}
$80\times 30 = \qty{2400}{N}$؛ وعند القدم $2400\times 5/45 =
\qty{270}{N}$.
\end{solution}

\begin{solution}{exo:b2:muscle-movement:9}
عند \qty{10}{Hz} تأتي الدفعات كل \qty{100}{ms}، بعد
زوال كالسيوم السابقة: فرعشات منفصلة، كل منها
تبدأ من حالة مرتخية جزئيًا، فتعطي تموجًا. وعند
\qty{50}{Hz} تأتي كل \qty{20}{ms}، قبل أن يُضخ الكالسيوم
بعيدًا: فيبقى الكالسيوم مرتفعًا وتندمج القوة. ويفوق
\omterm{def:b2:muscle-movement:motorunit}{الكزاز} \omterm{def:b2:muscle-movement:motorunit}{الرعشة} لأن دفقة كالسيوم واحدة أقصر من
أن تكفي الجسور المستعرضة لتمطيط العناصر المرنة المتوالية
وتوليد القوة الكاملة؛ أما الكالسيوم المستدام فيتيح لها ذلك.
\end{solution}

\begin{solution}{exo:b2:muscle-movement:10}
الشغل \qty{20}{kJ}، وطاقة ATP \qty{80}{kJ}، أي $1.6$ مول من ATP. وفي الثواني
\qty{5}{s} الأولى: $0.8$ مول من \omterm{prop:b2:muscle-movement:energy}{فوسفات الكرياتين}. والباقي $0.8$ مول من ATP
بالتحلل السكري: أي $0.4$ مول من الغلوكوز، و$0.8$ مول من اللاكتات --- وفي
\qty{40}{L} يكون ذلك \qty{20}{mmol/L}.
\end{solution}

\begin{solution}{exo:b2:muscle-movement:11}
قناة إطلاق محصورة: فلا كالسيوم ولا تقلص --- أي شلل.
ومضخة محصورة: فيبقى الكالسيوم مرتفعًا ولا تستطيع العضلة الارتخاء --- أي
تقفّع. وتروبونين غير حساس: فلا يتزحزح التروبوميوزين، ولا
تقلص. وبلا كالسيوم في الدم: يتقلص العضل الهيكلي
تقلصًا سويًا (فكالسيومه داخلي)؛ أما \omterm{def:b2:heart:anatomy}{القلب}، الذي يحتاج
كالسيومًا خارجيًا ليطلق إطلاقه، فيقف.
\end{solution}

\begin{solution}{exo:b2:muscle-movement:12}
يحوّل الرأس إلى شغل ما يصل إلى نصف الطاقة الحرة لجزيء ATP؛
والباقي حرارة، وخسارات المضخات، والجسور المستعرضة
التي ترتبط دون أن تجر عند السرعة، والعناصر المرنة
تنزل بالكل إلى الربع --- وهو مردود محرك بنزين
جيد. وطريقتا تدريج القوة تتيحان معًا تحكمًا دقيقًا عند الجهد
المنخفض (بإضافة الوحدات الصغيرة واحدة واحدة، وكل منها تنطلق أسرع أو
أبطأ) ومجالًا واسعًا (بألف ضعف، من وحدة صغيرة واحدة عند
معدل \omterm{def:b2:muscle-movement:motorunit}{الرعشة} إلى كل الوحدات في \omterm{def:b2:muscle-movement:motorunit}{كزاز})؛ أما مفتاح واحد فكان يعطي
قوة واحدة أو لا شيء.
\end{solution}

\begin{solution}{pb:b2:muscle-movement:1}
\textbf{1.} $v = \sqrt{2\times 9.81\times 0.4} = \qty{2.8}{m/s}$؛
و$\tfrac12\times 70\times 2.8^{2} = \qty{275}{J}$.
\textbf{2.} الشغل $= 275 + 70\times 9.81\times 0.4 = \qty{550}{J}$ على مدى
\qty{0.4}{m}: أي \qty{1370}{N}، ضعف الوزن.
\textbf{3.} $F_0 = 500\times 30 = \qty{15000}{N}$؛ وعند القدم
\qty{1500}{N}.
\textbf{4.} \qty{4}{cm} في \qty{0.3}{s}: أي \qty{0.13}{m/s}، و$1.1$ طول في
الثانية، أي $0.14\,v_{\max}$.
\textbf{5.} $F/F_0 = 0.3125/(0.14 + 0.25) - 0.25 = 0.55$: أي \qty{8300}{N}
في العضلة، و\qty{830}{N} عند القدم --- وهو أقل بكثير من
\qty{1370}{N} اللازمة. فتقصير العضلة وحده لا يستطيع أداء هذه القفزة.
والحركة المعاكسة تمد الباقي: فإذ يهبط القافز، تخزن
الأوتار الممططة طاقة مرنة، وتحمّلها العضلة وهي تتقلص
تقلصًا شبه متساوي القياس بقوة عالية؛ ثم ترتد الأوتار
أسرع مما تستطيع أي ليفة أن تقصر.
\textbf{6.} $550/0.3 = \qty{1.8}{kW}$؛ أي \qty{92}{W/kg} من الباسطات.
\textbf{7.} $4/\num{48000}$ cm $= \qty{0.83}{\micro m}$ لكل \omterm{def:b2:cell-differentiation:fibre}{قسيم عضلي}،
و\qty{0.42}{\micro m} لكل نصف: أي 42 ضربة قدرها \qty{10}{nm}.
\textbf{8.} \qty{0.42}{\micro m} في \qty{0.3}{s}: أي \qty{1.4}{\micro m/s}؛
و140 ضربة في الثانية لو كان الارتباط مستمرًا.
\textbf{9.} $500\times 6\times 10^{10}\times \num{48000}\times 300 =
4.3\times 10^{20}$ رأس.
\textbf{10.} $0.55\times \num{15000}\times 0.04 = \qty{330}{J}$؛ وفي
الضربة الواحدة $0.4\times 8\times 10^{-20} = \qty{3.2e-20}{J}$؛ أي $1.0\times
10^{22}$ ضربة.
\textbf{11.} $1.0\times 10^{22}/4.3\times 10^{20} = 24$ ضربة لكل
رأس، من 42 متاحة: أي نحو \qty{60}{\%}. والاحتياطي
صغير --- فالعضلة تعمل قرب حدها، كما وجد
السؤال 5.
\textbf{12.} $1.0\times 10^{22}/6.0\times 10^{23} = 0.017$ مول،
أي \qty{8.7}{g}؛ وطاقته $0.017\times 50 = \qty{860}{J}$، صار منها
\qty{330}{J} شغلًا: أي \qty{39}{\%}.
\textbf{13.} أسرع من دفعة واحدة كل \qty{30}{ms}: أي فوق
\qty{33}{Hz}؛ ونحو 10 دفعات في \qty{0.3}{s}.
\textbf{14.} $9.9\times 10^{-6}\times 5\times 10^{-10}\times 6\times
10^{23} = 3\times 10^{9}$ شاردة حرة؛ و$1.5\times 10^{9}$ من ATP.
\textbf{15.} الليفات: $500/(\pi\times(3\times 10^{-3})^{2}) = 1.8\times
10^{7}$؛ وATP الجسور المستعرضة لكل ليفة لكل دفعة $1.0\times
10^{22}/(1.8\times 10^{7}\times 10) = 6\times 10^{13}$، أي أربعين ألف
ضعف لتقدير المضخة. فالكالسيوم الحر بقية صغيرة: إذ
تطلق الشبكة ما رتبته \qty{0.1}{mmol/L}، أي عشرة
أمثال الارتفاع الحر، ويرتبط جُلّه في الحال بالتروبونين
وبالدارئات --- وكله يجب أن يُضخ راجعًا.
\textbf{16.} لقفزة قصوى، تُجنَّد كل الوحدات تقريبًا في غضون بضع عشرات
من المليثوان، الصغيرة أولًا والكبيرة آخرًا؛ أما الخطوة
اللينة فتجنّد نحو عُشرها إلى خُمسها.
\textbf{17.} تُجنَّد الوحدات بترتيب الحجم: فعند الجهد المنخفض تضيف كل
وحدة جديدة زيادة صغيرة، وعند الجهد العالي تكون الوحدات الباقية
كبيرة وتضيف كل منها خطوة كبيرة.
\textbf{18.} ينتصف كمون الصفيحة الانتهائية في كل ليفة؛ فالليفات
التي يهبط كمونها دون العتبة لا تنطلق، فتصغر
\omterm{def:b2:muscle-movement:motorunit}{الرعشة} \omterm{def:b2:muscle-movement:motorunit}{والكزاز}، وتقصر القفزة.
\textbf{19.} $5\times 20 = \qty{0.1}{mol}$ من ATP: أي ما يكفي ست قفزات.
\textbf{20.} $20\times 20 = \qty{0.4}{mol}$ من \omterm{prop:b2:muscle-movement:energy}{فوسفات الكرياتين}:
أي نحو 23 قفزة.
\textbf{21.} عشر قفزات تستعمل $0.17$ مول، ولا تنفد بها بعدُ \omterm{prop:b2:muscle-movement:energy}{فوسفات
الكرياتين}؛ وبعد عشرين أو نحوها يتولى التحلل السكري، وتتراكم اللاكتات
والبروتونات.
\textbf{22.} \qty{1.8}{kW} من الشغل عند \qty{40}{\%} هي \qty{4.6}{kW} من
ATP، أي $0.09$ مول في الثانية، وتحتاج $0.015$ مول أكسجين في الثانية ---
أي \qty{0.34}{L/s}، و\qty{21}{L/min}، وهو عشرة أمثال ما يستطيع الدم إحضاره:
فالقفزة لاهوائية بالضرورة، والأكسجين يأتي بعدها.
\textbf{23.} $0.17$ مول من ATP تحتاج $0.029$ مول من الأكسجين،
أي \qty{0.64}{L}؛ وإذا سُدد على مدى دقيقة كان \qty{640}{mL/min}، أي مثلي
استهلاك الراحة ونصف مثل.
\textbf{24.} يُعاد فسفرة ATP من \omterm{prop:b2:muscle-movement:energy}{فوسفات الكرياتين} في غضون
مليثوان، فلا يكاد تركيزه يهبط؛ والتصلب يحتاج ATP
قرب الصفر، وهو ما لا تبلغه قط عضلة حية فيها أي احتياطي.
\textbf{25.} 42 ضربة متاحة لكل نصف \omterm{def:b2:cell-differentiation:fibre}{قسيم عضلي}، و24 لازمة؛
و$F/F_0 = 0.55$ عند سرعة القفزة؛ و$0.017$ مول، أي \qty{8.7}{g}، من ATP.
\end{solution}
